(1) 选定一致的点号及正参考方向，有 $\lambda_1=L_1I_1+MI_2$、$\lambda_2=MI_1+L_2I_2$。被动电感的端电压（电流流入正端）为 $V_1=d\lambda_1/dt=L_1\dot I_1+M\dot I_2$、$V_2=d\lambda_2/dt=M\dot I_1+L_2\dot I_2$；若讨论反向的感应电动势，其符号为 $-d\lambda/dt$。
(2) 电源瞬时向两电感供能为 $P=V_1I_1+V_2I_2=dU/dt$。沿电流从零逐步建立的路径积分得到 $U=\tfrac12L_1I_1^2+MI_1I_2+\tfrac12L_2I_2^2$。将它配方为 $U=(L_1/2)(I_1+MI_2/L_1)^2+(L_2-M^2/L_1)I_2^2/2$，对任意电流无负磁能要求 $M^2\le L_1L_2$。
(3) 闭合超导第一回路无电阻且准静态变化使 $d\lambda_1/dt=0$，所以初末 $L_1I_1+MI_2=L_1I_1'$。因此 $I_1'=I_1+(M/L_1)I_2$。互感符号体现绕向，若 $M<0$，受迫调节方向相反；这里不允许擅自假设第一线圈电流保持不变。
(4) 把 $I_1=(\lambda_1-MI_2)/L_1$ 代入磁能，得到 $U=\lambda_1^2/(2L_1)+\tfrac12(L_2-M^2/L_1)I_2^2$。最终 $I_2'=0$，所以释放的可回收能量为 $\Delta U=U_{\rm before}-U_{\rm after}=\tfrac12L_2(1-k^2)I_2^2\ge0$，其中 $k=M/\sqrt{L_1L_2}$、$|k|\le1$；不能把它理解为在任何电路控制模式下均自动释放同样能量。
